%%
%% Author: shoupu_wan
%% 7/14/18
%%

\chapter{Mindsets for Dynamic Programming} \label{ch:dp}

Beginners like me would be intimidated by dynamic programming.
It seems what DP requires is a power of magic that creates multiple things at once and I just don't
have enough hands to handle all that at once.
DP needs a split of me doing one part of the thing and a split of me doing tightly coupled another
while the two splits coordinate closely.

\section{What types of problem may be solved by dynamic programming?}
\label{sec:types-of-problem-solved-by-dp}
For a problem to be solvable by dynamic programming, one must define
subproblems with the original problem as one of the subproblems (usually the largest, most
significant subproblem of all).
Also, the subproblems must be such that a directed acyclic graph (DAG) may be constructed among
them with dependencies among them.
Since the whole collection of subproblems is a DAG, it's necessary that some subproblems do not
depend on others and their solution is known right away, which we call the base subproblems.
Since a DAG establishes a partial ordering among its vertices, by traversing the graph starting
from the base subproblems following this partial ordering, we can solve all the subproblems  all
the way through to the largest subproblem, which is our original problem.
This is because all their dependencies has been resolved when we encounter them along the way.

There are several style of traversing---recursive or iterative.
In terms of the  auxiliary data structures to record the intermediate results, we may use
associative maps or hash sets.
We may use pre-allocated arrays and matrices which are generally faster to access.

So the complete thought path from recursion to dynamic programming is
\begin{itemize}
    \item Implement recursive solution
    \item Add memoir or cache
    \item Turn everything upside down and ``Start with the end'': build solutions bottom-up
    \item Vectorized your solution using arrays and matrix.
\end{itemize}

\section{Algorithm Refactor}\label{sec:algorithmRefactor}

\begin{defn}
    When a formula is too complicated, it may be refactored into several formalae, with auxiliary
    formulae which may make formulation of the problem simpler and the concept easier to understand.
    We call this strategy \emph{formula refactoring}.
\end{defn}

\begin{defn}
    When a dynamic programming formulation for a problem becomes too complicated, one may
    introduce a vector variable for each subproblem.
    Each component of the vector variable has a dedicated and meaningful purpose.
    The subproblem is to be solved to obtain the vector variable.
\end{defn}

\section{Two ways of dynamic programming}\label{sec:twoWaysOfDynamicProgramming}

The problem is asking for dynamic programming.
There are generally two ways of attacking dynamic programming.
Think of what's the optimal solution for subproblems and then grow the subproblem into the
full-sized problem;
think of what's the suboptimal solution for the same problem with certain
constraints and then seek to remove the constraints to obtain the optimal solution.
The first approach is 'problem-oriented' and the second is 'solution-oriented'.

Often, one can get clues by trying one of several ways.
One way is to seek restricted subproblems.
These subproblems would otherwise be just the original one.
But they are bounded by a single restriction.

Another way is to seek a partial solution to the original problem and then gradually build up the
solution to be a full one.

In both of these two types of reasoning, one may assume that all the smaller subproblems
have already had solutions (or subsolutions have been built) and one is free to use any subset of
them as aid to build the next subsolution.

\section{Dimensionality Boosting}\label{sec:dimensionality-boosting}

Dimensionality boosting is the opposite as dimensionality reduction.
In doing so, we increase the complexity of solution space with the hope to find an enclosing
solution space for the original problem.

When one runs out of wit in the thought process, one needs a way to `jump out of the
box' by adding dimensions.
It's like with a flat earth one can not go very far in explaining the astronomical phenomena
around.
But when one accepts that the earth is not flat, everything starts to make sense.
Then it is required for one to think in higher dimensional space, because the function
representing the recursive relationship becomes multivariate.
Sometimes even the function codomain turns into an array, a matrix, or a tensor.
Accordingly, the solution has to be stored in a collection, such as an array, a set or a map.

For example, in the problem ``Find the number of paths to escape $2$D grid'', because there is no
restriction as to the direction of move, a previously trailed path can be repeated.
This calls for the addition of the number of steps as an extra dimension when considering the
subproblem.

\begin{clisting}
    \lstinputlisting[label={lst:path_to_out_of_bound},
    caption={Dimensionality boosting to search for a solution space to contain the problem and
    its solution.}]
    {path_to_out_of_bound.py}
\end{clisting}



You are a salesperson selling stuff to households along a street.
Each house has a certain budget to buy your stuff.
All houses at this place are arranged in a circle.
That means the first house is the neighbor of the last one.
The neighborhood has a rule banning salesperson from selling stuff to adjacent houses.

Given a list of non-negative integers representing the budget of each house, determine
the maximum sales you can generate without violating the neighborhood rule.

Example 1:

Input: $[2,3,2]$
Output: 3
Explanation: You cannot sell to house 1 (budget = 2) and then to house 3 (budget = 2),
because they are adjacent houses.
Example 2:

Input: $[1,2,3,1]$
Output: 4
Explanation: Sell to house 1 (budget = 1) and then sell to house 3 (budget = 3).
Total sales $= 1 + 3 = 4.$


\begin{clisting}
    \lstinputlisting[label={lst:sales_tour},
    caption={Dimensionality boosting to search for a solution space to contain the problem and
    its solution.}]
    {sales_tour.py}
\end{clisting}


\section{Implementation}\label{sec:dp-implementation}

As to implementation, there are also two categories.
Category one is recursion with memoir.
Memoir can be anything an associative map, an array, etc.
Another is iterative (usually with the help of an array or matrix or higher dimensional
contiguous data structure).
Recursion with memoir is easier as long as one has a recursive relationship in mind.
Translating recursive relationship to iterative solution is a major hurdle to pass for dynamic
programming.
It requires abstract dimensional thinking, linear algebra, group theory and symmetry, topology,
set theory, graph theory, function and mapping.
Techniques that may be involved are index shifting.
The process of translating recursive relationship to iterative solution itself is iterative---it
is an iterative refining process.

In ~\autoref{ch:loop-thinning}, we persuade users away from using nested loops.
But for scientific computation, we mentioned that that direction can be relaxed depending on the
situation.
We hold similar stance here for dynamic programming as in dynamic programming we usually work
with high dimensional data structures and performance requirement is high.
It should be considered as highly specialized programming.
Nested loops are a good choice with careful consideration of alternatives.

\begin{clisting}
    \lstinputlisting[label={lst:new_21_game},
    caption={Translation from recursive relation to iterative dynamic programming code}]
    {new_21_game.py}
\end{clisting}

\begin{clisting}
    \lstinputlisting[label={lst:path_to_out_of_bound},
    caption={Choice of data structure in dynamic programming}]
    {path_to_out_of_bound.py}
\end{clisting}

\section{Recursive relation}\label{sec:recursive-relation}

Given a programming problem, how to define the subproblems and figure out the recursive relation?
The following problem helps illustrate this.

Greedy algorithm would certainly be suboptimal as given in this example.
\begin{enumerate}
    \item $[1, 1]$
    \item $[3, 3]$
    \item $[3, 3]$
\end{enumerate}

At first glance, the subproblems and their relations are unclear.
When I first attempted this problem, I was puzzled by how to `undo' the previous optimal solution
by regroup the books into rows as it seems necessary that an optimally packed subsolution does
not necessarily extend into the next subsolution by adding a book to it.

But as soon as I realized that I have all the subsolutions \emph{before} the current subproblem
available to build the next subsolution.
After one figure out a structure among the subproblems, it's like one come out of a darkness
into sunlight after a long tunnel.
\begin{clisting}
    \lstinputlisting[label={lst:fill_bookshelf},
    caption={Choice of data structure in dynamic programming}]
    {fill_bookshelf.py}
\end{clisting}

There are a row of houses, each house can be painted with three colors red, blue and green
\label{Evernote:Paint House (I and II)}
\begin{itemize}
    \item Problem-oriented\label{Evernote:Split "bothearthandsaturnspin" into words word break wordbreak}
    \item \item Solution-oriented
\end{itemize}

\section{Case Study--Tiling a Rectangle} \label{sec:tiling-rectangle}
\label{Evernote:2x1 brick paving problem tiling}
2x1 Bricks covering a rectangular area.

I want to illustrate the intimate interaction between implementation and algorithm.
One can have an tentative algorithm or even a proven working one.
As one implements, one may gain new insights and adjust the algorithm for better implementations.
Sometimes this can lead to unexpected results.
This is one of the examples.

\section{Two-sum Problem}

There is a greedy algorithm for solving this problems. But why does it work at all?

\section{Knapsack Problem}

The Knapsack problem is usually formulated as follows:
\begin{quote}
    \url{https://en.wikipedia.org/wiki/Knapsack_problem}
    Given a set of items, each with a weight and a value, determine the number of each item to
    include in a collection so that the total weight is less than or equal to a given limit and the
    total value is as large as possible.
\end{quote}

This problem is $NP$-hard.
But its limited form-integer weight and value-has a pseudo-polynomial solution.
A classical way to solve this problem is dynamic programming.
The question is ``Why the solution works at all?''
Why order of picking items affects neither the performance nor the result?
